% !TeX root = ../papers/nc-p3.tex
\section{階数15の縦切片 $x=1$ の完全分類}
\label{sec:rank-fifteen-vertical-all-levels}
\begin{theorem}\label{thm:rank-fifteen-vertical-all-levels}
剰余標数3の付値で $\bar x=1$, $\nu(y)=0$ となる標数零の通常な階数15窓は、
全て標準Koszul左右AS正則で、各始点で四中心四次式上自由階数256である。
\end{theorem}
\begin{proof}
第\ref{sec:rank-fifteen-vertical-unit-levels}節の楕円曲線 $W^2=t-t^3$ 上で、
$i^2=-1$ とし、係数を $1,i,W,iW$ の四座標で保持する。
\path{computations/verify_rank_fifteen_vertical_all_levels.py} は
$\mathbb F_3(t,i,W)$ 上で未切断経路を五次まで構成し、文字0と1の中心方程式を解く。
前節と同じ四本の中心を課して四始点の商を十三次まで消去する。
各文字で必要な本数の主元を選び、逆元を取るたびに分子の $\mathbb F_3(t)$ への
ノルムを記録する。入力の逆元から帰納すると、これらのノルムの零点を除けば
全ての除算が可能である。得られた平方自由な積 $F(t)$ の次数は 204451 である。

$N=t(1-t^2)$, $D=(1+t^2)^2$ と置けば、$v=N/D$,
$\bar y=P/Q$ において $P=(ND^2-N^3)^2$, $Q=D^6$ である。
この次数24の写像で $P,Q$ は互いに素で、$P'Q-PQ'$ の零点は $PQ=0$ 上に限る。
従って有限の非零な $\bar y$ の各ファイバーは相異なる24点を持つ。
$t\mapsto-t,1/t,(1+t)/(1-t)$ のいずれでも $P/Q$ は変わらない。
$F$ から $PQ=0$ の点を除き、これら三変換の引き戻しとの最大公約元を反復すると、
安定した式の次数は 784 である。全ての表示が $F=0$ に入るファイバーはこの式に含まれる。
残る各既約因子で $P/Q$ の最小多項式 $g$ を求め、同じ $g$ に対応する因子の積が
$Q^{\deg g}g(P/Q)$ のモニック化と等しいかを厳密に比較する。

全表示が除かれる剰余軌道は、$Y+1$, $Y-1$, $Y^3+Y^2+2Y+1$,
$Y^6+Y^3+2Y^2+2Y+1$ の四つに限る。一次と三次は前節の証明で認証済みである。
六次では $\mathbb F_{3^6}$ の表示で最小多項式と七つの倍率座標を直接検査する。
未切断経路の長さ0から5、中心の文字別次元 $(3,2,0,0)$、
四中心商の四始点の長さ13までの次元を取り直す。有限体の共役で六根が尽くされる。

従って各剰余値では少なくとも一表示で全ての主元が取れる。
未切断経路・中心の一般の下界、中心商の上界と持ち上げは
定理\ref{thm:rank-fifteen-unit-residues}の通りであり、剰余値が超越的なときも
非零な $F$ の零点を外れるため同じ判定が成り立つ。
\end{proof}

\begin{corollary}
標数零の代数閉体において、通常な階数15の縦切片 $x=1$ の全点が上記の正則性と自由性をもつ。
添字保持同型類は $t\in\mathbb G_m\setminus\{1,-1\}$ で一意に表され、
水準 $y=-(t-t^{-1})^2/108$ ごとに $y=1/27$ では二類、それ以外の $y\ne0$ では四類である。
全類の最小周期は2である。
\end{corollary}
\begin{proof}
係数体に3進付値を延長する。$\nu(y)>0$ なら
定理\ref{thm:rank-fifteen-unit-pole-depth}で $J=0$、$\nu(y)<0$ なら極の定理、
$\nu(y)=0$ なら本節の定理を適用する。同型座標と一位置移動
$t\mapsto-t$ は系\ref{cor:rank-fifteen-vertical-one-rational}による。
\end{proof}
